\documentclass[a4j,12pt,dvipdfmx]{jsreport}

% 数式
\usepackage{amsmath,amsfonts}
\usepackage{bm}

% 画像
\usepackage[dvipdfmx,top=35truemm,bottom=30truemm,left=30truemm,right=30truemm]{geometry}
\usepackage[dvipdfmx]{graphicx}
% biblatex
\usepackage[style=numeric, sorting=none, maxbibnames=1000]{biblatex}
\addbibresource{graduation_thesis.bib}

% itembox
\usepackage{ascmac}
% breakitemboxのためのpackage
\usepackage{itembkbx}
% \usepackage{tcolorbox} 画像が見えなくなってしまう原因
% \usepackage{framed} 枠

% フォント
\usepackage[ipaex]{pxchfon}
\listfiles
\usepackage{float}
\usepackage{multicol}
\usepackage{multirow}
\usepackage{paralist}
\usepackage{caption}
% \captionsetup{format=plain, justification=centering}
\usepackage{setspace}
\usepackage{algorithm}
\usepackage{algpseudocode}
\usepackage[dvipdfmx,hidelinks]{hyperref}
\usepackage{pxjahyper}
\usepackage{fancyhdr}
\usepackage{ltablex}
\usepackage{capt-of}
\usepackage{tabularx}
\usepackage[export]{adjustbox}
\usepackage{booktabs}
\usepackage{enumitem} % itemizeの行間 これは下になければならない
\usepackage{threeparttable} % 表の脚注
\usepackage{titlesec}
\usepackage{makecell}
\usepackage{forest}
\usepackage{fancybox}
% titlesec と \paragraph の相性問題を回避（run-in 見出しにする）
\titleformat{\paragraph}[runin]
  {\normalfont\normalsize\bfseries}
  {}{0pt}{} % ←ここには文字や \noindent 等を入れない
\titlespacing*{\paragraph}
  {0pt}{1.0ex}{0.8em}
\titleformat{\chapter}[hang]{\LARGE\bfseries}{\prechaptername\thechapter\postchaptername}{1em}{}

\setlength{\headheight}{16.5pt}
\captionsetup{format=hang}
\captionsetup[figure]{font=small}
\captionsetup[table]{font=small}
\begin{document}

% ▼▼▼ 仕掛け：main じゃないとき（単体のとき）だけヘッダー設定 ▼▼▼
\ifdefined\ismain\else
   \pagestyle{fancy}
   \lhead{\rightmark}
   \renewcommand{\chaptermark}[1]{\markboth{第\ \thechapter\ 章\ #1}{}}
   \setlength{\headheight}{16.5pt}
   \titleformat{\chapter}[hang]{\LARGE\bfseries}{\prechaptername\thechapter\postchaptername}{1em}{}
\fi
% ▲▲▲▲▲▲▲▲▲▲▲▲▲▲▲▲▲▲▲▲▲▲▲▲▲▲▲▲▲▲▲▲▲▲▲▲

\chapter*{研究業績一覧} \label{performance}
\addcontentsline{toc}{chapter}{研究業績一覧}

\section*{国内会議}

\begin{enumerate}[label=(\roman*)]
   \item \textbf{渡邉謙吾}，川村天, 小林伶央, 有井知真, 伊藤亮史, 栗原聡，"物語構造分析に基づくLLMを用いたインタラクティブストーリー生成システム"，第24回データ指向構成マイニングとシミュレーション研究会(SIG-DOCMAS)(人工知能学会合同研究会2023)，Nov.25，2023
   \item 安部玲央, 有井知真, 伊藤亮史, \textbf{渡邉謙吾}, 小林伶央, 川村天, 栗原聡，"アフォーダンスを活用するマルチエージェントプランニングにおけるタスク分解"，第24回 データ指向構成マイニングとシミュレーション研究会(SIG-DOCMAS), Nov.25，2023
\end{enumerate}

% ▼▼▼ 仕掛け：main じゃないとき（単体のとき）だけ参考文献を表示 ▼▼▼
\ifdefined\ismain\else
   \printbibliography[title=参考文献]
\fi
% ▲▲▲▲▲▲▲▲▲▲▲▲▲▲▲▲▲▲▲▲▲▲▲▲▲▲▲▲▲▲▲▲▲▲▲▲

\end{document}